\clearpage\thispagestyle{otFront}\phantomsection\label{references}\pdfbookmark[0]{Sources et édition}{bm-references}\otPageTitre{Sources et édition}{Références officielles · France}\fontsize{11}{15}\selectfont
\par\medskip\textbf{1. Arrêté et calendrier}\par « Programmes d’enseignement de français et de mathématiques du cycle des approfondissements (cycle 4) » : arrêté du 18 février 2026, Bulletin officiel n\textsuperscript{o} 10 du 5 mars 2026 (NOR MENE2602912A). Le nouveau programme s’applique en Quatrième à la rentrée 2027-2028.\par\href{https://www.education.gouv.fr/bo/2026/Hebdo10/MENE2602912A}{\textcolor{otGreenDark}{Consulter la source officielle}}\par
\par\medskip\textbf{2. Programme de mathématiques du cycle 4}\par Annexe 2 de l’arrêté, rubriques « Quatrième » : automatismes et objectifs d’apprentissage en nombres et calculs, espace et géométrie, données et probabilités, proportionnalité et fonctions, pensée informatique.\par\href{https://www.education.gouv.fr/sites/default/files/document/Annexe%202%20%E2%80%93%20Programme%20de%20math%C3%A9matiques%20pour%20le%20cycle%204-480716.pdf}{\textcolor{otGreenDark}{Consulter la source officielle}}\par
\par\medskip\textbf{3. Exemples pour la mise en œuvre}\par Éduscol, « Mathématiques. Exemples pour la mise en œuvre des programmes 2026. Cycle 4 » : exemples de réussite et conseils de mise en œuvre pour la Quatrième.\par\href{https://eduscol.education.gouv.fr/sites/default/files/document/exemple-mise-en-oeuvre-c4-mathematiques0pdf-126080.pdf}{\textcolor{otGreenDark}{Consulter la source officielle}}\par
\par\medskip\textbf{4. Ressources d’accompagnement}\par Éduscol : ressources d’accompagnement du programme de mathématiques au cycle 4.\par\href{https://eduscol.education.gouv.fr/5736/ressources-d-accompagnement-du-programme-de-mathematiques-au-cycle-4}{\textcolor{otGreenDark}{Consulter la source officielle}}\par
\par\bigskip Sources consultées le 23 octobre 2026. Les cours sont rédigés pour ce
livre ; les exercices, défis et bilans sont originaux, sans reproduction des
exemples officiels.\par\medskip
Périmètre de Quatrième : produits et quotients de relatifs ; opérations sur les
rationnels ; puissances d’exposant positif ; racine carrée et encadrement ;
distributivité simple et preuves littérales ; équations $ax+b=cx+d$ ; ratios,
pourcentages et grandeurs quotients ; programmes de calcul et dépendance ;
moyenne pondérée, médiane et étendue ; probabilités avec notations ensemblistes
et deux épreuves ; droites remarquables et logique ; translations ; droite des
milieux ; Pythagore, réciproque et contraposée ; cercle circonscrit au triangle
rectangle ; pyramides et cônes ; conditions et variable en programmation.
Les défis approfondissent ces outils sans ajouter d’exigence officielle.

\clearpage
